\begin{abstract}
Public bike-sharing systems require effective overnight repositioning to restore station inventory levels before the next operating period. This thesis studies static bike repositioning problems with user-dissatisfaction-based service measures, with emphasis on scalability, broken-bike handling, and on-site repair operations.

The first part develops a user-dissatisfaction-aware cluster-first, route-second decomposition framework for the multi-vehicle static bike repositioning problem. The framework incorporates the User Dissatisfaction Function into the clustering phase, evaluates candidate clusters through an Activated Transfer Selection procedure, and solves the resulting clustering problem using a Hybrid Genetic Search. The computational study assesses the proposed decomposition against exact optimization and baseline clustering methods.

The second part extends the static bike repositioning problem by considering broken bikes and repairers who perform repairs directly at stations. The model jointly represents truck routing, repairer routing, and station-level decisions for repairing, loading, and unloading bikes. A hybrid solution algorithm, HGSADC-SBC, is used to integrate route search with station-budget-constrained quantity assignment.

The third part combines these two directions. It develops an extended CFRS framework for the static bike repositioning problem with broken bikes and on-site repair. The framework accounts for the expected change in user dissatisfaction caused by repositioning, repair, and station-level inventory changes. Together, the three components of the thesis provide formulations and solution methods for service-quality-oriented static bike repositioning under increasingly rich operational settings.
\end{abstract}
